\section{Part X: Advanced Research and Innovation Decision Cases}

\subsection{71. The 5-page scientific investment case}

The most consequential research proposal should not begin with projected revenue. Page 1 states why the question matters and what the field already knows. Page 2 states the decisive experiments, independent baselines and potential sources of bias. Page 3 covers talent, rights, ethics and facilities. Page 4 contains full cost, funding restrictions, milestones and downside. Page 5 records what the institute will learn regardless of the outcome, commercialization routes, conflicts, and the specific decision being requested.

Require the scientific sponsor to distinguish known facts, plausible mechanisms and speculative claims. The capital committee should not approve a proposal whose scientific foundation it cannot explain in plain language after hearing an independent reviewer.

\subsection{72. Bayesian decision analysis for a research option}

[Model] Assume a potential technology works under realistic conditions with prior probability $0.25$. A small evaluation has sensitivity $0.80$ and specificity $0.90$. The unconditional probability of a positive diagnostic test is $0.25 \times 0.80 + 0.75 \times 0.10 = 0.275$. Following a positive result, Bayes' probability updates to:
\begin{equation}
P(\text{Works} \mid \text{Positive}) = \frac{0.20}{0.275} \approx 0.727
\end{equation}
Following a negative result, the posterior probability falls to:
\begin{equation}
\begin{aligned}
&P(\text{Works} \mid \text{Negative}) \\
&\quad = \frac{0.25 \times 0.20}{0.25 \times 0.20 + 0.75 \times 0.90} \approx 0.069
\end{aligned}
\end{equation}
This shows why a diagnostic experiment can change a decision. But the assumptions are invented and its sensitivity/specificity may be unknowable. A laboratory should never use a crisp Bayesian posterior to disguise fragile validation data. Vary the inputs and examine what conclusions survive.

\subsection{73. Value of information with an expensive downstream commitment}

[Model] Suppose successful development yields net present value ₹80 lakh after the follow-on investment, while failure produces -₹40 lakh. At $25\%$ prior success probability, unconditional investment has expected value:
\begin{equation}
0.25 \times 80 - 0.75 \times 40 = -\textnormal{₹}10\text{ lakh}
\end{equation}
so the rational simple decision is not to invest. If a perfect ₹3-lakh experiment revealed success or failure in advance, investing only on success produces:
\begin{equation}
0.25 \times 80 - 3 = \textnormal{₹}17\text{ lakh}
\end{equation}
expected incremental value over doing nothing. The theoretical value of perfect information before cost is ₹20 lakh; this study would be worth ₹17 lakh if truly perfect.

Real studies are imperfect and delayed, and the option may expire. An experiment's actual expected value requires a decision tree using its error rates and costs. The point is that a small research budget can earn its return by preventing a bad large bet.

\subsection{74. Research option decay and competitor timing}

Some discoveries have short commercial windows: a foundation-model provider may incorporate the feature, a patent may be filed elsewhere, an industry standard may be fixed, or a market may be captured. Others, such as basic mathematical knowledge, can remain useful for decades. A project should therefore record both the \textbf{time required to learn} and the \textbf{time until learning loses relevance}.

Do not equate speed with scientific validity. Accelerate decision-changing tests and dissemination when appropriate, while preserving the controls that make the evidence reliable. A fast unreliable result may destroy more option value than it creates.

\subsection{75. Portfolio concentration and correlated failure}

[Model] Imagine 4 AI application projects all depend on a single model provider's API and on 1 enterprise customer's permitted dataset. Their research questions differ, but a provider policy change or customer withdrawal can disable all 4 simultaneously. Counting 4 projects as independent diversification is false. A portfolio map must include shared model, dataset, partner, principal investigator, equipment and funding dependencies.

Mitigate concentration where worthwhile: provider-independent evaluation, redundant rights-cleared datasets, shared but portable infrastructure, alternate collaborators and contractual exit permissions. Diversification should not be so broad that it destroys specialist expertise.

\subsection{76. When research should stay inside the commercial company}

Keep a project inside a commercial operating company when its primary output improves a product or service, its IP and customer contracts are already governed there, and specialized academic or charitable financing is not required. This is common for engineering R\&D and process improvement. Track the actual cost and identify whether capitalization of development expenditure is permitted under applicable accounting standards instead of automatically expensing or capitalizing all work.

Move to a separate research entity only with a real legal, funding or independence rationale. Transferring assets between related entities can create tax, shareholder, contract and donor questions.

\subsection{77. When a separate nonprofit research center is defensible}

A separately governed mission entity may be appropriate if the research has public-benefit objectives, requires restricted grants or gifts, employs independent scholars, and needs publication or academic autonomy distinct from the commercial parent. It must be genuinely compliant and cannot be used to subsidize founder-owned products without lawful, properly approved transactions.

Before incorporation, secure a credible scientific program, potential funders, qualified governance, financial controls and an interentity rights map. An empty ``foundation'' without scientific capacity creates administrative expense and the illusion of progress.

\subsection{78. The patent-versus-publication decision in numbers}

[Model] Assume a patent program costs ₹6 lakh initially and ₹2 lakh annually for 5 years, for ₹16 lakh nominal outlay before enforcement. A license opportunity that produces ₹20 lakh expected gross royalty receipts may not justify protection once legal fees, commercialization probability, delay and alternative open dissemination are considered. Conversely, a patent enabling a regulated company to invest in development could be highly valuable even when direct institutional royalties are small.

Use an invention decision memo. The relevant question is not whether a patent can be obtained, but whether protection makes a socially and economically desirable translation more likely.

\subsection{79. Commercialization partners may capture most value}

A small lab can create a technically valuable method yet lack manufacturing, sales, certification, service and distribution. Bargaining power may therefore sit with the licensee. Negotiate milestones and retained use rights, and compare multiple plausible partners where possible. Do not inflate a valuation based only on the downstream market's size.

Teece's complementary-assets view is helpful: invention and profitable commercialization depend on different capabilities. A lab should know which ones it lacks before choosing exclusive licensing or a spinout.

\subsection{80. Founder-led research versus independent scientific leadership}

A founder with a statistics degree, engineering background and AI experience can reasonably originate hypotheses, build prototypes and commission critiques in those domains. That does not automatically make them the correct scientific leader for every future research discipline. Significant research agendas need appropriately qualified investigators, methods and independent review.

The founder's comparative advantage may eventually be capital allocation, recruiting, problem selection and institution-building rather than serving as principal investigator for every project. This is a progression in responsibility, not a loss of scientific identity.

\subsection{81. Measuring a researcher's potential without overvaluing publications}

Evaluate paper quality and specific contribution; replication and methodology; ability to revise beliefs; mentoring; software or experimental craft; ethical judgment; and ability to formulate useful questions. Publication count alone is contaminated by field norms, team size, incentives and institutional privilege. A candidate may be excellent at research engineering yet publish little.

Use a work sample: ask a candidate to critique a flawed study, propose a decisive test under resource constraints, and revise the plan after contradictory evidence. Compensate substantive work where appropriate and protect applicants' confidential ideas.

\subsection{82. How to terminate an unsuccessful research project fairly}

Stopping a project should preserve safety, data rights, participants, outstanding contracts, research records and staff dignity. Give investigators a chance to present new evidence; decide using pre-identified standards and current opportunity costs; offer publication of useful negative findings when lawful; redeploy talent if appropriate. Do not equate a stopped project with an incompetent scientist.

Conversely, do not preserve a weak project solely because the institution has already purchased equipment or wants a positive headline. Sunk costs are not future benefits.

\subsection{83. Science and customer discovery can reinforce each other without becoming the same thing}

For translational work, interview potential users early to understand adoption constraints, existing substitutes, procurement and economics. Such evidence helps choose applications, but customer demand does not determine whether a scientific claim is true. Research design may reveal that a highly demanded result is not technically possible, or that an unfashionable negative finding is important.

Keep 1 document for scientific evidence and another for buyer evidence, linking them only where causal support exists.

\subsection{84. Pricing sponsored research and consulting differently}

A fixed-scope technical implementation can be priced on deliverables, risk and staff time. Research sponsorship often purchases a process of inquiry with uncertainty about results; milestones should therefore include auditable work and reporting rather than a guaranteed discovery. Pricing must cover actual indirect costs and the opportunity cost of granting exclusive IP or publication restrictions.

A sponsor may reasonably ask for budget controls and confidentiality. It should not be promised a successful scientific conclusion as a contractual deliverable.

\subsection{85. A research center's break-even may be strategically misleading}

A mature sponsored research unit might receive enough grants and contracts to cover its direct costs but still require a large institutional subsidy for faculty time, facilities and administration. Alternatively, a research unit generating excellent public scholarship may intentionally run at an approved deficit. ``Break-even'' is not the universal goal; \textbf{approved and financeable net subsidy per credible research outcome} is often the better question.

Report sponsor recovery and total economic cost separately. Reconcile to unrestricted reserves and multi-year commitments.

\subsection{86. A research facility's capital-replacement trap}

Assume ₹1 crore of lab equipment with a 5-year practical replacement cycle. Even ignoring maintenance and financing, maintaining equivalent capability requires planning for an average ₹20 lakh annual replacement need. If the lab charges clients only enough to cover consumables and salaries, reported contribution may be positive while the facility steadily decapitalizes.

Use a capital replacement schedule. Some equipment may last longer or become obsolete earlier, so a straight-line estimate is only a planning baseline.

\subsection{87. Choosing a center of excellence partner}

Before signing a college or enterprise ``center of excellence'' agreement, identify who owns instruments, who schedules access, who trains technicians, who provides working capital, who selects researchers, who can publish, who bears incidents, and who maintains equipment after sponsorship. A room with a logo is not a research center.

Demand a credible 1st-year research program with named questions, people, cash, rights and external evaluation. If the counterpart mainly wants marketing symbolism, decline or call it a training lab.

\subsection{88. India-specific institutional grant eligibility}

Official ANRF calls can be highly specific about eligible principal investigators and host institutions; BIRAC opportunities differ by technology and scheme. A founder should maintain a matrix: call name, eligibility, deadline, host requirements, permitted costs, matching funds, indirect recovery, IP/publication and reporting. Do not presume a new private commercial company qualifies for all grants advertised to universities or nationally recognized research organizations.

For each attractive scheme, obtain written confirmation or a reliable interpretation of eligibility before allocating substantial proposal-writing effort. This avoids expensive rejection for structural reasons.

\subsection{89. International research collaboration due diligence}

Cross-border projects introduce research ethics, export controls in sensitive fields, sanctions, data transfers, intellectual property, publication norms, grant restrictions, immigration, taxes and currency exposure. Choose a lead institution and clear responsibility for approvals, records and incident reporting. A visiting professor's enthusiasm cannot bind their university without authorization.

Start with low-risk seminars, reproducible open work or limited data-free collaborations if complex permissions are unresolved. Escalate to shared data, facilities and funding only after governance is settled.

\subsection{90. From 1 research question to a durable research program}

A single project may succeed because of 1 researcher. A program survives successive projects through an articulated scientific agenda, high-quality mentorship, method libraries, accessible infrastructure, external critique, reproducibility practices and funding relationships. Allocate some effort to \textbf{program assets} rather than only project outputs: common benchmarks, safe datasets, training materials and review systems.

The founder should test whether the next investigator can extend the work without starting entirely from zero. That is scientific institutionalization.

\subsection{91. The 24-month research decision map}

Months 1-3: choose and review 1 question. Months 4-6: reproducible baseline and 1st negative/positive evidence. Months 7-12: replication or partnership, contingent on proof. Months 13-18: sustained external review and suitable funding proposal if eligible. Months 19-24: decide whether the institute has a real program warranting a permanent research lead. These are planning windows, not deadlines for scientific discovery.

If the organization has no demonstrated commercial surplus, researchers or research rights, pause the capital-intensive branch and preserve public technical scholarship as the low-cost research option.

\subsection{92. A founder's research red-team session}

Every 6 months ask an outside specialist to argue that the institute's best research narrative is wrong: baselines are weak, selection effects drive outcomes, claims are not novel, methods cannot replicate, the partner owns the valuable data, or an AI provider will commoditize the result. Reward the finding of a real flaw.

A research reputation is strengthened when its internal reviewers can prevent overstatement before external exposure.

\subsection{93. The ethical commercialization boundary}

An educational institution should not extract student equity through grading power, use unconsented learner data to train commercial models, impose indefinite publication suppression for sponsor revenue, or misrepresent faculty research as institution-owned success without proper attribution. A future successful spinout is not worth a present breach of trust.

Make commercialization optional for students, govern it through independent terms, and ensure academic learning and credentials remain available without giving the institute private venture upside.

\subsection{94. Research decision principle}

Fund the next experiment because it can \textbf{change what the institution rationally believes or does}, not because it gives the institute a research label. Build a lab because a specific body of scientific work needs that lab and can finance its responsibilities. Create a venture when a distinct market opportunity and team exist. Preserve the capacity to stop, publish accurately and redirect capital when evidence changes.

A research enterprise becomes institutionally valuable when it can produce both meaningful results and reliable corrections to its own errors.
